% =====================================================================
%  Chapter 17 — The Relic Density Constraint
% =====================================================================

\chapter{The Relic Density Constraint}
\label{ch:relic}

The hardest constraint in dark matter physics is not a null result from a
detector. It is the number 0.12. Every theory that claims to describe dark
matter must produce it, and the ways in which theories produce it are more
varied, and more revealing, than the ways in which they evade the exclusion
curves.

\section{The statement of the constraint}The dark matter density today is

\[
\Omega_c h^2 = 0.1200 \pm 0.0012 .
\]

This is a two per cent measurement, and it is a constraint on the entire
history of the universe: every mechanism that produced dark matter, every
mechanism that destroyed it, and every process that redistributed it between
species, has to be consistent with the number.

In a thermal equilibrium history, the density today is set by the rate at
which the species decoupled. The relic abundance of a species with
annihilation cross section $\langle\sigma v\rangle$ and mass $m$ is

\[
\frac{\Omega_\chi h^2}{0.12} \;\approx\;
\frac{3\times 10^{-26}\ \cmthree\s}{\langle\sigma v\rangle}
\left(\frac{m}{100\ \GeV}\right)^{1.2},
\]

up to the standard model-independent and scheme-dependent factors. The
practical content of this formula is the scaling, not the prefactor: the
density is inversely proportional to the annihilation rate and nearly
independent of the mass.

\section{Why this is a miracle, and why that word is doing work}For the
weak-scale thermal relic, the formula gives the observed density with no free
parameter. That coincidence is the WIMP miracle of Chapter~\ref{ch:wimp}, and
it has three logical readings, which the field has not agreed on.

The \emph{anthropic} reading is that it is a coincidence of the Standard
Model parameters and that there is no deeper explanation. The
\emph{fine-tuned} reading is that the scale is fixed by other considerations
--- the electroweak scale is itself unexplained --- and the coincidence is a
consequence of a deeper connection. The \emph{correct} reading is that the
picture is wrong somewhere else, and the coincidence is a clue to where.

The honest position in 2026 is that the coincidence remains unexplained and
the WIMP's failure to appear has removed the candidate that would have made
it meaningful. This is the sharpest form of the origins problem: the one
mechanism that explained the density for free is the one that the
experiments have eliminated.

\section{What the constraint rules out}Because the density is known precisely, any
mechanism that produces too much of it is excluded, and this is a much
broader exclusion than the laboratory provides.

It rules out thermal relics that would have decayed to a smaller abundance
today, which is the case for any species with a decay channel and a
lifetime shorter than the age of the universe. It rules out any
species whose abundance was increased at any point after the freeze-out
without a compensating mechanism. It rules out the simplest versions of
asymmetric dark matter when the asymmetry is larger than the observed
density. And it constrains any alternative cosmology: a period of late
matter domination, a period of extra entropy, or a non-standard expansion
history all change the relic density, and the constraint on the cosmology
from the dark matter abundance is correspondingly strong.

The last of these deserves emphasis, because it cuts against the popular
intuition that a new dark sector automatically means new flexibility. A
dark sector with its own thermal history is not a free parameter: it must
still produce 0.12, and the mechanisms by which it does so are exactly the
mechanisms that the CMB and the baryon acoustic oscillations constrain.

\section{What the constraint does not rule out}It does not rule out a mechanism
that produces the density through a parameter that is itself unexplained
and untestable: the axion's initial misalignment angle, the primordial
curvature power spectrum that fixes the primordial black hole mass
function, the scale of a light field's production.

This is the asymmetry at the centre of the field. Mechanisms that predict
the density are constrained by the density, because a prediction can be
wrong. Mechanisms that parameterise the density are not constrained by it,
because a parameter cannot be wrong. The field's discovery strategy has to
account for this, and Chapter~\ref{ch:portfolio} does.

\begin{keytakeaways}
\item The relic density is known to two per cent and constrains every
production mechanism, every post-freeze-out history, and every alternative
cosmology.
\item The WIMP miracle is the only mechanism that explained the density with
no free parameter, and its candidate has been excluded.
\item Mechanisms that predict the density are falsifiable by it; mechanisms
that parameterise it are not, and the discovery strategy must treat them
differently.
\item The abundance constraint on any alternative cosmology is strong enough
that a new dark sector is not a free ticket.
\end{keytakeaways}
